% =========================================================================
% SECCIÓN 4: TAXONOMÍA Y MATRIZ COMPARATIVA DEL ESTADO DEL ARTE
% =========================================================================

\section{Taxonomía y Matriz Comparativa del Estado del Arte}
\label{sec:taxonomia_comparativa}

\subsection{Taxonomía de Enfoques de Asistencia Lingüística Basados en IA}
\noindent [Descripción de las dimensiones taxonómicas principales: (1) Paradigma de inferencia (Nube vs. Local/Offline); (2) Estrategia de conocimiento (Paramétrico puro, RAG, Híbrido); (3) Grado de adaptación lingüística (Zero-shot, In-context learning, LoRA/QLoRA, Preentrenamiento continuo); (4) Enfoque pedagógico (Formativo guiado, Consulta directa, Evaluación).]

\subsection{Matriz Comparativa de Investigaciones Previas}
\noindent La Tabla~\ref{tab:matriz_sota_comparativa} sintetiza y compara los trabajos más representativos identificados en la revisión frente a la propuesta de esta investigación.

\begin{table}[t]
\centering
\caption{Matriz comparativa de investigaciones del estado del arte.}
\label{tab:matriz_sota_comparativa}
\resizebox{\textwidth}{!}{%
\begin{tabular}{lllllll}
\toprule
\textbf{Estudio / Autores} & \textbf{Modelo Base} & \textbf{Estrategia} & \textbf{Idioma / Variedad} & \textbf{Métricas Clave} & \textbf{Despliegue} & \textbf{Limitaciones Principales} \\
\midrule
{}[Estudio 1, Año] & [LLaMA-2 7B] & [Zero-shot Prompting] & [Quechua Chanka] & [BLEU: X, chrF: Y] & [Servidor Cloud] & [Alta latencia, dependencia de internet] \\
{}[Estudio 2, Año] & [Qwen-1.5 1.8B] & [RAG canónico] & [Español general] & [ROUGE-L: X, Exact Match] & [Local CPU] & [Sin soporte para lenguas originarias] \\
{}[Estudio 3, Año] & [Gemma 2B] & [LoRA Fine-tuning] & [Lenguas amerindias] & [Perplejidad, F1] & [GPU 8\,GB] & [Alucinación frecuente en datos técnicos] \\
{}[Estudio 4, Año] & [Phi-2 2.7B] & [RAG + Cuantización 4b] & [Guaraní] & [Precision@k, Recall@k] & [Edge device] & [Falta de validación pedagógica escolar] \\
{}[Estudio 5, Año] & [Mistral 7B] & [QLoRA + RAG] & [Aymara] & [chrF++, ROUGE-1] & [GPU 16\,GB] & [Inviable para hardware modesto offline] \\
\midrule
\textbf{Nuestra Propuesta} & \textbf{[Qwen2.5-3B-Q4]} & \textbf{[RAG + QLoRA]} & \textbf{[Quechua Collao]} & \textbf{[chrF++, Latencia, RAGAS]} & \textbf{[Local Offline ($\le$4GB)]} & \textbf{[Supera brecha offline y morfológica]} \\
\bottomrule
\end{tabular}%
}
\end{table}

\subsection{Análisis Crítico de Trade-offs Tecnológicos}

\subsubsection{Compromiso entre latencia y fidelidad pedagógica}
\noindent [Análisis cuantitativo de tiempos de respuesta por token vs. coherencia conceptual en explicaciones gramaticales.]

\subsubsection{Compromiso entre huella de memoria (VRAM) y capacidad de razonamiento}
\noindent [Evaluación del punto de inflexión entre cuantización extrema (2--3 bits) y degradación en el seguimiento de instrucciones complejas.]

\subsubsection{Compromiso entre recuperación densa vs. léxica en morfología rica}
\noindent [Discusión sobre la efectividad de embeddings semánticos frente a algoritmos de correspondencia léxica (BM25) en palabras aglutinadas.]

\subsection{Síntesis de Hallazgos por Pregunta de Investigación}

\subsubsection{Síntesis para RQ1 (SLMs y Despliegue)}
\noindent [Conclusiones del análisis sobre viabilidad de modelos de 1.5B a 3B en GPUs de 4\,GB de VRAM.]

\subsubsection{Síntesis para RQ2 (RAG para Lenguas Originarias)}
\noindent [Conclusiones sobre técnicas de recuperación y segmentación morfológica.]

\subsubsection{Síntesis para RQ3 (Adaptación PEFT/QLoRA)}
\noindent [Conclusiones sobre el afinamiento eficiente sin pérdida de capacidades generales.]

\subsubsection{Síntesis para RQ4 (Evaluación y Métricas)}
\noindent [Conclusiones sobre la combinación necesaria de métricas automáticas (chrF++) y rúbricas pedagógicas humanas.]
